\section{Native transcription of the finite-truncation interface and mixed-channel closure file (source pp.~232--280)}
This appendix replaces the former archival Appendix~F source-page block with a native REV\TeX{} presentation of the theorem-facing finite-truncation interface and mixed-channel closure material. In the architecture of the core paper, the principal outputs here are the canonical finite-truncation SFTE/remainder package, the finite mixed bounded-family packaging node, and the finite-cap mixed-correlator closure theorem. The same source-page range also contains downstream Route~1 reformulation material; to preserve source continuity without disturbing proof-home routing, those items are recorded here in a short auxiliary subsection.

\subsection{Finite-cap renormalized local field algebra and intrinsic remainder control}
The first part of the source range turns the canonical finite-truncation construction into an explicit finite-cap local-field algebra with a common invariant core.

\begin{theorem}[Source Theorem~F.243: dyadic construction of the renormalized local field algebra up to $\Delta_{\max}$]
Assume the weak-coupling curvature-control and weak-window package used in Lane~A. Fix $\Delta_{\max}\ge 0$ and the canonical blockwise basis of the exact-dimension quotient formalism. Then there exist:
\begin{enumerate}[label=(\roman*),leftmargin=2.2em]
\item a graded triangular matrix function $Z^{(\le \Delta_{\max})}(t)$ on $(0,t_\ast)$,
\item for each test function $\varphi\in C_c^\infty(\mathbb R^4)$ constants $C_{\varphi,\Delta_{\max}}<\infty$ and $\eta>0$,
\end{enumerate}
such that the dyadic renormalized multiplet obeys a summable increment bound and therefore converges in $L^2$ as the dyadic flow time tends to zero. The limit defines a renormalized local field multiplet $O_{\mathrm{ren}}^{(\le \Delta_{\max})}(\varphi)$ on an explicit common dense invariant core of bounded flowed positive-time classes.
\end{theorem}

\begin{lemma}[Source Lemma~F.244: explicit finite-cap common invariant core]
Fix $\Delta_{\max}\ge 0$ and let $\mathcal D_{\mathrm{flow}}^{(\le \Delta_{\max})}$ be the span of bounded flowed positive-time classes in the finite-cap OS Hilbert space. Then:
\begin{enumerate}[label=(\roman*),leftmargin=2.2em]
\item $\mathcal D_{\mathrm{flow}}^{(\le \Delta_{\max})}$ is dense and invariant under the bounded dyadic approximants;
\item for each canonical field component and each bounded flowed positive-time factor, the dyadic approximants form a Cauchy sequence in the OS Hilbert norm and in every fixed finite graph norm relevant to the polynomial field algebra;
\item consequently the renormalized field components are closable, preserve the common invariant core, and define a well-posed finite polynomial field algebra there.
\end{enumerate}
\end{lemma}

\begin{corollary}[Source Corollary~F.245: canonical finite-truncation SFTE and intrinsic remainder control]
Fix $\Delta_{\max}\ge 0$ and the canonical blockwise basis. Then the canonical flowed lifts and the renormalized local fields produced by Theorem~F.243 satisfy the Lane~A small-flow-time expansion on that finite truncation,
\[
Y_t^{(\le \Delta_{\max})}(f)
=
C^{(\le \Delta_{\max})}(t;\mu)
O_{\mathrm{ren}}^{(\le \Delta_{\max})}(f;\mu)
+
R_t^{(\le \Delta_{\max})}(f;\mu),
\]
with a block-lower-triangular coefficient matrix and a remainder whose matrix elements against every fixed finite insertion family decay after multiplication by the same finite-cap inverse-control function.
\end{corollary}

\proofsynopsis{The source proof packages the canonical one-shell transport, imposes the graded triangular renormalization conditions recursively in the grading order, proves summability of the dyadic shell increments, and upgrades the resulting $L^2$ control to graph-norm control by the insertion-to-OS transfer theorem. The corollary is the exact finite-cap SFTE/remainder package later consumed by the mixed-channel closure theorem.}

\subsection{Auxiliary downstream Route~1 reformulation retained from the same source range}
The source-page range 276--279 also contains a downstream Route~1 reformulation block. These results are not core proof homes for Companion~II, but they are retained here in native form to preserve continuity with the March~16,~2026 source file.

\begin{theorem}[Source Theorem~F.326: secondary Route~1 dense-core bridge via QE3 and OS spectral calculus]
Let $\omega^{(u_{\mathrm{ren}})}$ be the sharp-local continuum state furnished by the Lane~A closure chain and let $(\mathcal H,\Omega,H)$ be its OS reconstruction. Then the dense positive-time core generated by bounded local cylinder observables is OS-dense and satisfies Euclidean exponential clustering with the same gap rate $m_\ast$ already obtained on the primary Route~1 spine.
\end{theorem}

\begin{theorem}[Source Theorem~F.327: identification of the Route~1 bounded state and the Lane~A sharp-local state]
For the tuned dyadic sequence attached to $u_{\mathrm{ren}}$, the bounded positive-time state and its dyadic OS matrix-element limits agree with the restriction of the full sharp-local continuum state to the positive-time algebra. Consequently every Route~1 estimate first proved for the bounded dyadic OS matrix elements transfers to the same sharp-local state used in the main theorem.
\end{theorem}

\begin{corollary}[Source Corollary~F.329: noncollapse of the physical mass scale from one-shot entrance]
Along the tuned ultraviolet sequence associated with $u_{\mathrm{ren}}$, there exists a uniform strictly positive lower bound on the physical mass scale. In the source architecture this is the theorem-facing noncollapse consequence of the one-shot entrance theorem together with the exact entrance-scale covariance transfer.
\end{corollary}

\begin{remark}
These Route~1 reformulation items are included here only because they occupy the same source-page block as Corollary~F.330A and Theorem~F.331. The stable proof home for the core paper's Route~1 theorems remains Companion~I.
\end{remark}

\subsection{Finite mixed bounded-family packaging and the sole mixed-correlator closure node}
The last part of the source range contains the actual finite-cap mixed-channel seam.

\begin{corollary}[Source Corollary~F.330A: finite mixed bounded-family packaging on a finite truncation]
Fix $u_{\mathrm{ren}}\in(0,u_\ast]$ and $\Delta_{\max}\ge 0$, and let
\[
B_{\mathrm{cap}}^{(\le \Delta_{\max})}:=\ast\bigl(A_{\mathrm{flow}}^{(\le \Delta_{\max})},A_+\bigr)
\]
be the unital bounded capwise algebra generated by the capwise flowed subalgebra and the bounded positive-time algebra. For every finite family $W\subset B_{\mathrm{cap}}^{(\le \Delta_{\max})}$, the tuned Wilson expectations of the words generated by $W$ converge jointly along the dyadic trajectory, and therefore define a unique packaged state
\[
\omega^{(u_{\mathrm{ren}})}_{W,\mathrm{cap},\mathrm{bnd}}
\]
on the generated finite algebra. These packaged states are positive, normalized, compatible under enlargement of the family $W$, and inherit Euclidean translation covariance on translated finite families.
\end{corollary}

\begin{theorem}[Source Theorem~F.331: finite-cap mixed-correlator extension from finite-family packaged bounded capwise states]
Fix $u_{\mathrm{ren}}\in(0,u_\ast]$ and $\Delta_{\max}\ge 0$. Let
\[
B_{\mathrm{cap}}^{(\le \Delta_{\max})}:=\ast\bigl(A_{\mathrm{flow}}^{(\le \Delta_{\max})},A_+\bigr)
\]
and let $\widehat A_{\mathrm{loc}}^{(\le \Delta_{\max})}$ be the unital $\ast$-algebra generated by $B_{\mathrm{cap}}^{(\le \Delta_{\max})}$ and the formal centered capwise local symbols. Then, after first fixing a finite family $W\subset B_{\mathrm{cap}}^{(\le \Delta_{\max})}$ large enough to contain every bounded word appearing in the formula under consideration, the packaged state from Corollary~F.330A extends to a mixed-correlator functional on the algebra generated by those words and the centered local symbols, with the following theorem-facing consequences:
\begin{enumerate}[label=(\roman*),leftmargin=2.2em]
\item replacing the designated local slot by its bounded inverse-flow representative changes every mixed correlator by $o(1)$ as the flow parameters tend to zero;
\item the resulting value is independent of the admissible finite family $W$ chosen to package the bounded capwise state;
\item the positive-time mixed-channel matrix elements are compatible with the same bounded positive-time state that later enters the Lane~A finite-cap sharp-local extension theorem.
\end{enumerate}
This theorem is the first mixed-correlator closure node on the Lane~A spine and the sole theorem-level mixed-channel closure node cited by the core paper.
\end{theorem}

\proofsynopsis{The source proof first packages every finite bounded word family into a positive compatible capwise state, then proves that the inverse-flow approximants are Cauchy in every fixed mixed channel against those packaged states. Because compatibility under enlargement of the finite family is built into Corollary~F.330A, the limiting mixed-correlator functional is independent of the chosen auxiliary bounded family. This is precisely what turns the finite-cap seam into a theorem-facing closure node rather than a diffuse bookkeeping argument.}

\begin{remark}
Within the architecture of the core paper, the role split must remain literal: Corollary~F.330A supplies the finite mixed bounded-family packaging node, while Theorem~F.331 is the sole mixed-correlator closure theorem at finite cap. The theorem is fixed-cap and fixed-family in that precise sense: the auxiliary bounded family $W$ is frozen before the closure statement is read, and no inductive-union or full-union quantitative uniformity is asserted here. The inductive-union passage and bounded-base cyclicity occur only later in the Lane~A chain.
\end{remark}
